% Standalone mathematical component for the parent report.
% Requires amsmath, amssymb, amsthm; define theorem/lemma if absent.
\section{An explicit weak-arithmetic transport theorem}
\label{sec:ipra-transport}

The following result separates a new machine transformation from the
still-unproved consistency calibration of its ancestor. It is a theorem
about intuitionistic primitive recursive arithmetic, not an inference from
testing finitely many executions.

\subsection{Exact hypotheses}

Work in the usual language of $\mathrm{iPRA}$, with a function symbol and
defining equations for every primitive recursive function, and induction
for quantifier-free formulas. A predicate below is represented by a
primitive recursive characteristic function. All the displayed hypotheses
are required to be \emph{theorems of $\mathrm{iPRA}$}, uniformly in their
free variables. Truth for standard inputs alone does not suffice.

Let $F$ and $G$ be primitive recursive total step functions on codes of
source and target configurations respectively. Their domains may differ;
both are coded by natural numbers. Let $h_F,h_G$ be their Boolean halting
tests. Totalize a halted computation by leaving its configuration fixed.
It is enough to require the following weaker absorbing-flag law:
\[
 h_G(y)=1\ \longrightarrow\ h_G(G(y))=1.                 \tag{A}
\]
There is no additional physical working state in this semantic
totalization: a designated halt label acquires an absorbing interpretation
only in the arithmetic run function.

Let $c_0,b_0$ be the source and target initial configuration codes; let $P$
be a primitive recursive source-frame validity predicate; let $E$ encode
each valid source configuration as a target boundary configuration; and
let $\tau$ be a primitive recursive segment-duration function. Assume
\begin{align}
 &P(c_0)=1,                                                 \tag{B}\label{eq:sim-initial}\\
 &P(c)=1\ \longrightarrow\ P(F(c))=1,                       \tag{C}\label{eq:sim-closed}\\
 &P(c)=1\ \longrightarrow\ 1\leq\tau(c),                 \tag{D}\label{eq:sim-positive}\\
 &P(c)=1\ \longrightarrow\ G^{\tau(c)}(E(c))=E(F(c)),     \tag{E}\label{eq:sim-exact}\\
 &P(c)=1\ \longrightarrow\ h_G(E(c))=h_F(c).              \tag{F}\label{eq:sim-halt}
\end{align}
Here $F^n$ and $G^n$ denote iteration, which is primitive recursive.
Finally suppose that, for a specified natural number $d$,
\[
 G^d(b_0)=E(c_0).                                          \tag{G}\label{eq:sim-startup}
\]
A verified fixed blank-tape initialization supplies (G). An explicit
primitive recursive parameter-dependent startup duration can also be
used; the fixed form is enough for a blank-start machine.

These are syntactic proof obligations. In particular, (E) is an exact
endpoint equality, not ``the target either halts or eventually reaches
the endpoint.'' The latter disjunction would not supply the reverse
implication below. The totality and primitive recursiveness of $\tau$ are
part of the hypotheses; they cannot be inferred merely because the
intended source macro is expected to terminate.

\begin{theorem}[Positive primitive-recursive simulation]
Under (A)--(G), $\mathrm{iPRA}$ proves
\[
 \bigl(\exists n\ h_F(F^n(c_0))=1\bigr)
 \ \longleftrightarrow\
 \bigl(\exists t\ h_G(G^t(b_0))=1\bigr).
\]
Consequently $\mathrm{iPRA}$ proves the equivalence of the two
nonhalting sentences. The proof uses no excluded middle for arbitrary
formulas, no Markov principle, and no consistency assumption.
\end{theorem}

\begin{proof}
Define the primitive recursive cumulative duration
\[
 T(0)=0,\qquad
 T(n+1)=T(n)+\tau(F^n(c_0)).
\]
Quantifier-free induction using (B),(C) proves $P(F^n(c_0))=1$.
Quantifier-free induction using (D) then proves $T(n)\geq n$.
The iteration identity
\[
 G^{u+v}(y)=G^v(G^u(y))
\]
is itself proved by quantifier-free induction on $v$. Applying it at
each successor stage with (E), a further quantifier-free induction gives
\[
 G^{T(n)}(E(c_0))=E(F^n(c_0)).
\]
Together with (G) this is
\[
 G^{d+T(n)}(b_0)=E(F^n(c_0)).                            \tag{H}
\]
Every induction predicate just used is an equality, an inequality, or a
finite conjunction of such primitive recursive assertions.

If $n$ witnesses source halting, (H) and (F) show that
$d+T(n)$ witnesses target halting. This is an explicit witness
transformation, primitive recursive in $n$.

For the converse, suppose $t$ witnesses target halting. Quantifier-free
induction from (A) proves
\[
 h_G(G^t(b_0))=1\ \longrightarrow\
 h_G(G^{t+r}(b_0))=1
\]
for every $r$. Since $d+T(t)\geq t$, substitute
$r=d+T(t)\mathbin{\dot{-}}t$ and use the elementary truncated-subtraction
identity to obtain $h_G(G^{d+T(t)}(b_0))=1$.
By (H) and (F), $h_F(F^t(c_0))=1$. Thus the identity function
$t\mapsto t$ supplies a source-halting witness. This also covers a
purported halt during startup; no least-boundary search is required.

The two positive existential implications prove the claimed equivalence
intuitionistically. If $H_F\leftrightarrow H_G$, intuitionistic logic
alone gives $\neg H_F\leftrightarrow\neg H_G$.
\end{proof}

\subsection{Allowing zero-length administrative segments}

The positivity assumption can be replaced by an explicit cofinality
certificate. Keep (A)--(C) and (E)--(G), permit $\tau(c)=0$, and supply a
primitive recursive $N$ together with an $\mathrm{iPRA}$ proof of
\[
 d+T(N(t))\geq t\quad\hbox{for every }t.                 \tag{I}
\]
The same proof works, with reverse witness transformation $t\mapsto N(t)$.
For example, if each group of at most $k>0$ source steps consumes at least
one target step, a proved bound $T(kn)\geq n$ supplies $N(t)=kt$.
The bound and its proof are required; merely knowing that every standard
target time eventually lies before some boundary is not a certificate in
the required weak arithmetic.

Zero-duration local simulation without cofinality is insufficient. Take
$F(n)=n$, $h_F(n)=0$, $E(n)=0$, $\tau(n)=0$, and target states $0,1$
with $G(0)=G(1)=1$, $h_G(0)=0$, $h_G(1)=1$. Set both initial codes to
zero. The exact endpoint equation holds for every source configuration,
and the halting flags agree at every encoded boundary, but the target
halts after one step while the source never halts. The missing condition
is precisely cofinality of the selected target boundaries.

\subsection{How a literal machine transformation can meet the hypotheses}

Use a canonical finite-support tape encoding, for example a state and two
binary stacks around the head. A literal transition, an iterated run,
frame construction, and tape equality are then primitive recursive. This
avoids an unbounded extensional quantifier over tape positions in (E).
The initial all-zero tape has a fixed code.

For a symbolic runtime certificate, every accepted symbolic step must
have its own arithmetic derivation. A literal table row is a closed
equation. Traversing an affine unary run is proved by quantifier-free
induction on its length. A register-file excursion needs a proved
primitive recursive scan bound and the corresponding exact endpoint
equation, uniformly in the whole finite file. A finite certificate
concatenates those witnesses and their durations. This construction
produces an actual term $\tau(c)$ rather than an unbounded search for a
successful duration.

The parent machine's frame invariant must include enough information for
all called subroutines: header and delimiter positions, the complete
register/spare list, the represented source control, and any guard used by
a history-sensitive fold. If the proof depends on a two-cell context,
that context belongs in $P$ or in the source configuration. Graph
reachability of a control is not a substitute for the claimed context
invariant.

A finite certificate-checker run can verify the candidate-specific closed
facts, but it does not establish the checker rules' uniform soundness in
$\mathrm{iPRA}$. Conversely, once those rule proofs, the frame closure,
and the primitive recursive duration bounds are supplied, no separate
analysis of arbitrary interior target configurations is necessary. The
exact endpoint law, the absorbing halt flag, and cofinality already
reflect every finite target halt.

\subsection{What this theorem does and does not settle}

Suppose $M$ is an existing machine and $M'$ a transformed one. A completed
instantiation gives the genuine relative result
\[
 \mathrm{iPRA}\vdash
 \bigl(\neg\operatorname{Halt}(M)\leftrightarrow
       \neg\operatorname{Halt}(M')\bigr).
\]
To conclude
\[
 \mathrm{iPRA}\vdash
 \bigl(\neg\operatorname{Halt}(M')\leftrightarrow
       \operatorname{Con}(\mathrm{PA})\bigr),
\]
one still needs an unconditional $\mathrm{iPRA}$ calibration of $M$, or
an independent calibration of $M'$. An ordinary PA calibration of $M$
does not discharge this premise. The distinction persists even if all
new machine transformations are proved in $\mathrm{iPRA}$.

